\mainchapterstyle
\chapter{تجزیه و تحلیل داده‌ها}
\pagestyle{mainheader}


\section{مقدمه}
ارائه‌ی داده‌ها، نتایج و تحلیل و تفسیر اولیه آنها در این فصل ارائه می‌شود. در ارائه‌ی نتایج با توجه به
راهنمای کلی نگارش فصل‌ها، تا حد امکان ترکیبی از نمودار و جدول استفاده شود. با توجه به حجم و ماهیت
تحقیق و با صلاحدید استاد راهنما، این فصل می‌تواند تحت عنوانی دیگر بیاید. در صورتی که حجم داده‌ها زیاد
باشد، بهتر است به صورت نمودار یا در قالب ضمیمه ارائه نشده و فقط نمونه‌ها در متن آورده شود. این فصل
باید به سوالات تحقیق عطف به یافته‌های محقق پاسخ داده شود. اگر تحقیق دارای آزمون فرض باشد، پذیرش
یا عدم پذیرش فرضیه‌ها در این فصل گزارش می‌شود. این فصل حدود 40 صفحه است.

\section{محتوا}
در این بخش به سوالات تحقیق براساس داده‌ها و یافته‌های محقق، پاسخ داده می‌شود.داده‌ها با فرمت
مناسبی ارائه می شوند. مدل‌ها اجرا شده و نتیجه آن مشخص می‌شود.

\section{اعتبارسنجی}
از طریق مقایسه نتایج با نتایج کارهای دیگران، استفاده از روش‌های تحلیل پایائی\LTRfootnote{Reliable} و اعتبار\LTRfootnote{Validity} ، نظرگیری از
خبرگان\LTRfootnote{Expert Judgment} انجام می شود.
